\subsection{Final source projection: eight columns and no private bound}
\label{sec:projected-direct}

The source projection in this subsection is explicitly credited to the
incoming \emph{ConZFC Direct Projection 7617} package, whose preprocessing
proof is preserved as \texttt{sources/INCOMING\_PROJECTION\_PROOF.md}.
The new result here is its composition with the sharper carrier,
parametric positive-power graphs, dyadic majorants, quotient bound $x$,
and sign-normalized integer compilation.  We do not claim that eliminating
these three label columns was first discovered in this construction.

Each rule can be parameterized by just $a,b,c,d$:
\[
\begin{array}{c|c|c|c}
 \text{rule}&p&q&v\\ \hline
 \text{truth}&0&0&0\\
 \text{eight ordinary axioms}&0&0&f_j(a,b,c,d)\\
 \text{bundled axiom}&0&0&f_{\rm bundle}(a,b,c,d)\\
 \text{modus ponens}&a&I(a,b)&b\\
 \text{generalization}&b&0&A(a,b)
\end{array}
\]
The bundled row additionally requires
$((c-a)(c-b))^2=(d+1)^2$.  No other row has a guard.  These substitutions
are exactly the pinned C222 rules, including the independent
parameter and premise of generalization.

For a chosen branch $j$, substitute its labels $(p_j,q_j,v_j)$ into the
accumulator residual
\[
 L_j=S'(1+v_jM)-S(1+p_jM)(1+q_jM).
\]
Let $C_j=L_j^2$ except at the bundle, where
$C_j=L_j^2+g^2$.  The new local polynomial is
\[
 P_8=\prod_{j=0}^{11}C_j
 +\bigl(x-S-(1+(i+1)y)u\bigr)^2
 +\bigl(x-S'-(1+(i+2)y)u'\bigr)^2. \tag{12}
\]
Its zero set is exactly one valid branch-specific update together with
the two common beta reconstructions.  As all branch factors are
nonnegative, the product implements the finite disjunction without
selecting a rule by an extra coordinate.  This is existential projection
of the three deterministic labels, not an assertion that the polynomial
in (12) is algebraically identical to the eleven-coordinate source.

\paragraph{Padding chooses the parameter bound before the carrier.}
Given a finite proof-tree expansion, choose $N\ge1$ at least its length
and at least every schema parameter in the expansion.  For modus ponens,
this may include a very large formula code; no practical bound on the
resulting $N$ is claimed.  Pad to exactly $N$ rows using truth, whose
labels are $(0,0,0)$ and whose update leaves $S$ unchanged.  Thus all four
parameters may be bounded by $N$ without changing refutability.

Every branch label is a nonnegative-coefficient polynomial in those four
parameters.  Apply the structural exponent certificate to its literal
constructor DAG with radius $N+4$.  The largest certified exponent is
3,445,903, so the fixed polynomial
\[
 m=(N+4)^{2^{22}} \tag{13}
\]
bounds every label.  It is computed by exactly 22 squarings, not by a
variable-exponent graph and not by a duplicate evaluation of the bundled
axiom.  It satisfies $m\ge1$.  Choose $M=(m+1)^N$ and $M_0=M-1$.
The carrier proof already given applies verbatim to these labels.  The
choice order is proof, padded $N$, then $m,M$, then the accumulator
sequence, then $y,x$.  There is no self-referential bound.

The eight local coordinates are
\[
 (a,b,c,d,S,S',u,u'),\qquad
 \text{bounds }(N,N,N,N,y-1,y-1,x,x). \tag{14}
\]
The use of $x$ for the two quotient bounds is an additional simplification:
canonical beta reconstructions have
$x=S+Du$ with $S\ge0,D\ge1$, hence $u\le x$, and similarly for $u'$.
No new private bound $U$ is needed.  Conversely, the state bounds in (14)
make those reconstructions canonical, just as in the eleven-column proof.
The endpoints and parameterization $y=M_0+2+y_0$ are unchanged.
The four outer natural coordinates are $N_0,x,y_0,M_0$, with $N=N_0+1$.

\paragraph{A second dyadic radius proves the shared CRT bound.}
Split the actual signed circuit for (12) as $P_8=E_8-F_8$ and put
\[
 R=4+N+x+y+M_0,\qquad H=R^{2^{23}}. \tag{15}
\]
All inputs at the upper substitution of (14) are at most $R$; for the two
states we safely use $y$ instead of $y-1$.  The two root exponents from
the structural certificate are both 6,893,248, below $2^{23}$.
Therefore $H$ bounds both source evaluations and $N$.  It also dominates
every binomial cutoff $N+1,y,x+1$.  The same CRT theorem, with direct
factorial cancellation, applies.  Only 23 successive squarings compute
(15).  Notice that $m$ need not occur in this source radius: after label
projection the local source computes its labels from the four bounded
parameters and uses only $M_0$ as its carrier input.

The final eight-column arithmetic uses exactly 40 positive-power graphs,
11 binomial graphs and two factorial graphs.  The count is transparent:
24 powers for the eight coordinate binomials, five for each of the two
factorials, five other powers in the modulus product, and one carrier
power.  In this accounting the five other product powers are its two
explicit powers and the three powers of its extra binomial; the factorial
at $N$ has already been counted.  No arithmetic macro, variable-length
product or variable exponent remains after flattening.

\subsection{Existential sign normalization for the final integer equation}

The incoming signed-quartic package supplies the sign-orbit method
and its two pivots per positive-power graph; its proof is preserved as
\texttt{sources/INCOMING\_SIGN\_NORMALIZATION\_PROOF.md}.  The final
application here is checked against the new projected parametric system,
not assumed merely because it worked for a predecessor.

For a family of residuals $r_j$, a diagonal sign action
$x_i\mapsto(-1)^{\epsilon_i}x_i$ preserves the common zero set whenever
all monomials of each row have equal parity against $\epsilon$.
Then $r_j$ is changed to either $r_j$ or $-r_j$.
For each of the 40 power calls there are two independent actions:
\begin{itemize}
\item Flip its auxiliary $Z$ and the product wire $WZ$.  The only other
      use of that wire is through its square in the last norm.
\item Flip its Pell coordinate $U$ and its two unrestricted quotients
      $q_B,q_S$.  The products $Uq_B,Uq_S$ are unchanged, while the norm
      uses $U$ only through $U^2$.
\end{itemize}
The actions leave every other semantic coordinate unchanged.  They
commute, and do not change any retained sign guard.  At any integer zero,
apply a pivot's action exactly when that pivot is negative.  This produces
another zero in which all 80 pivots are nonnegative.  This is an
existence-preserving normalization, not a claim that the original zero
already had those signs.

The released normalization certificate specifies all flipped raw and
exported variables and the parity of every residual under each action.
The separate independent checker verifies these finite identities; the
argument above explains their universal meaning.  Gate values involving
signed quotients remain allowed to be negative.  No signs of $X,S,R_0$
are discarded by an unsupported symmetry argument.

The mixed-domain template has 1,937 variables and 1,660 quadratic rows.
It has 161 unrestricted congruence quotients: four per power call and
one final CRT quotient.  There are 529 other semantic coordinates.
Before normalization, exact polynomial sign closure requires 494 seeds
and infers 35 more signs.  The 80 pivots are among these seeds and can
instead be normalized.  Thus only
\[
 414+80+35=529
\]
semantic signs need to be guarded, normalized, or inferred, respectively.
The 414 three-triangular-number certificates add $3\cdot414=1242$
integer variables and 414 quadratic rows.  Consequently the final
integer counts are
\[
 1937+1242=3179,\qquad 1660+414=2074.
\]
The resulting ordinary quartic $Q_8$ has 23,157 monomials and coefficient
height 63,923,766.  The natural companion has 1,859 variables, 1,581 rows,
12,275 monomials and height 63,924,066.  Both heights are below $2^{26}$.
The literal files are \texttt{projected/expanded\_polynomial.json} and
\texttt{projected\_integer/expanded\_polynomial.json}.

\paragraph{Complete equivalence and comparisons.}
A zero of the integer polynomial first yields an integer zero of the
mixed template with the 414 guarded signs.  Normalize the 80 pivots and
replay the 35 sign implications.  All nonquotient arithmetic inputs are
then natural.  The parametric Pell certificates establish every required
power; formulas (7)--(9) establish every other arithmetic macro; the CRT
certificate recovers all eight-column rows; their parameter and label
bounds license the sharp carrier and hence a C222 refutation.  The pinned
calibration gives a ZFC refutation.  Conversely a finite refutation gives
the padded source in (12)--(14), all arithmetic certificates, natural
power witnesses, and the 414 triangular representations, hence an integer
zero.  Therefore
\[
 \operatorname{Con}(\mathrm{ZFC})
 \quad\Longleftrightarrow\quad
 \neg\exists z\in\mathbb Z^{3179}\;Q_8(z)=0.
\]
All steps are ordinary mathematics formalizable in ZFC.  Finite artifact
checks do not replace the stated unbounded arithmetic and calibration
theorems, and no new proof-assistant certification or proof of consistency
is claimed.

The late incoming direct integer alternatives have 3,776 variables,
20,194 monomials and height 47,959; and 3,671 variables, 26,271 monomials
and height 3,835,369,234.  The new 3,179-variable equation improves the
latter on variables, monomials and coefficient height, but does not
improve the former's monomial count or its much smaller coefficients.
It should not be described as dominating every current direct equation.
The 2,560-natural-variable sparse companion likewise remains relevant
because its 11,074 monomials are fewer than the projected natural
companion's 12,275.  These are explicit tradeoffs, not an absolute
minimality claim.
